\documentclass[CJK,13pt]{beamer}
\usepackage{CJKutf8}
\usepackage{beamerthemesplit}
\usetheme{Malmoe}
\useoutertheme[footline=authortitle]{miniframes}
\usepackage{amsmath}
\usepackage{amssymb}
\usepackage{graphicx}
\usepackage{eufrak}
\usepackage{color}
\usepackage{slashed}
\usepackage{simplewick}
\usepackage{tikz}
\usepackage{tcolorbox}
\usepackage{ulem}
\graphicspath{{../figures/}}
%%figures
\def\lfig#1#2{\includegraphics[width=#1 in]{#2}}
\def\tfig#1#2{\includegraphics[height=#1 in]{#2}}
\def\addfig#1#2{\begin{center}\includegraphics[width=#1 in]{#2}\end{center}}
\def\question{\includegraphics[width=0.3in]{why.jpg}\,}
\def\answer{\includegraphics[width=0.3in]{baozou_haha.png}\,}
\def\wulian{\includegraphics[width=0.18in]{emoji_wulian.jpg}}
\def\bigwulian{\includegraphics[width=0.35in]{emoji_wulian.jpg}}
\def\bye{\includegraphics[width=0.18in]{emoji_bye.jpg}}
\def\bigbye{\includegraphics[width=0.35in]{emoji_bye.jpg}}
\def\huaixiao{\includegraphics[width=0.18in]{emoji_huaixiao.jpg}}
\def\bighuaixiao{\includegraphics[width=0.35in]{emoji_huaixiao.jpg}}
\def\jianxiao{\includegraphics[width=0.18in]{emoji_jianxiao.jpg}}
\def\bigjianxiao{\includegraphics[width=0.35in]{emoji_jianxiao.jpg}}
\def\haoqi{\includegraphics[width=0.18in]{emoji_haoqi.jpg}}
%% colors
\def\blacktext#1{{\color{black}#1}}
\def\bluetext#1{{\color{blue}#1}}
\def\redtext#1{{\color{red}#1}}
\def\darkbluetext#1{{\color[rgb]{0,0.2,0.6}#1}}
\def\skybluetext#1{{\color[rgb]{0.2,0.7,1.}#1}}
\def\cyantext#1{{\color[rgb]{0.,0.5,0.5}#1}}
\def\greentext#1{{\color[rgb]{0,0.7,0.1}#1}}
\def\darkgray{\color[rgb]{0.2,0.2,0.2}}
\def\lightgray{\color[rgb]{0.6,0.6,0.6}}
\def\gray{\color[rgb]{0.4,0.4,0.4}}
\def\blue{\color{blue}}
\def\red{\color{red}}
\def\orange{\color[rgb]{1.,0.8,0.}}
\def\green{\color{green}}
\def\darkgreen{\color[rgb]{0,0.4,0.1}}
\def\darkblue{\color[rgb]{0,0.2,0.6}}
\def\skyblue{\color[rgb]{0.2,0.7,1.}}
%%control
\def\diag{\mathrm{diag}\,}
\def\heaviside{\,\mathrm{h}\,}
\def\bral{\left(\begin{array}{l}}
\def\brar{\end{array}\right)}
\def\brall{\left(\begin{array}{ll}}
\def\brarr{\end{array}\right)}
\def\bralll{\left(\begin{array}{lll}}
\def\brarrr{\end{array}\right)}
\def\branchl{\left\{\begin{array}{l}}
\def\branchr{\end{array}\right.}
\def\branchll{\left\{\begin{array}{ll}}
\def\branchrr{\end{array}\right.}
\def\branchlll{\left\{\begin{array}{lll}}
\def\branchrrr{\end{array}\right.}
\def\sfgamma#1{\,\Gamma\left( #1 \right)\,}
\def\be{\begin{equation}}
\def\ee{\nonumber\end{equation}}
\def\bea{\begin{eqnarray}}
\def\eea{\nonumber\end{eqnarray}}
\def\bch{\begin{CJK}{UTF8}{gbsn}}
\def\ech{\end{CJK}}
\def\bitem{\begin{itemize}}
\def\eitem{\end{itemize}}
\def\bcenter{\begin{center}}
\def\ecenter{\end{center}}
\def\bex{\begin{minipage}{0.2\textwidth}\includegraphics[width=0.6in]{jugelizi.png}\end{minipage}\begin{minipage}{0.76\textwidth}}
\def\eex{\end{minipage}}
\def\chtitle#1{\frametitle{\bch#1\ech}}
\def\bmat#1{\left(\begin{array}{#1}}
\def\emat{\end{array}\right)}
\def\bcase#1{\left\{\begin{array}{#1}}
\def\ecase{\end{array}\right.}
\def\bmini#1{\begin{minipage}{#1\textwidth}}
\def\emini{\end{minipage}}
\def\tbox#1{\begin{tcolorbox}#1\end{tcolorbox}}
\def\pfrac#1#2#3{\left(\frac{\partial #1}{\partial #2}\right)_{#3}}
\def\res#1#2{\,\mathrm{res}\,#1\left(#2\right)\,}
\def\newt#1#2{\left(\begin{array}{c}#1\\ #2\end{array}\right)}
\def\reof#1{\,\mathrm{Re}{\left(#1\right)\,}}
\def\imof#1{\,\mathrm{Im}{\left(#1\right)\,}}
\def\Arg#1{\,\mathrm{Arg}\,#1\,}
%%symbols
\def\intfull{\int_{-\infty}^\infty}
\def\inthalf{\int_0^\infty}
\def\bropt{\,(\ \ \ )}
\def\sech{\mathrm{sech}\,}
\def\csch{\mathrm{csch}\,}
\def\asinh{\mathrm{asinh}\,}
\def\acosh{\mathrm{acosh}\,}
\def\sone{$\star$}
\def\stwo{$\star\star$}
\def\sthree{$\star\star\star$}
\def\sfour{$\star\star\star\star$}
\def\sfive{$\star\star\star\star\star$}
\def\rint{{\int_\leftrightarrow}}
\def\roint{{\oint_\leftrightarrow}}
\def\stdHf{{\textit{\r H}_f}}
\def\deltaH{{\Delta \textit{\r H}}}
\def\ii{{\dot{\imath}}}
\def\skipline{{\vskip0.1in}}
\def\skiplines{{\vskip0.2in}}
\def\lagr{{\mathcal{L}}}
\def\hamil{{\mathcal{H}}}
\def\vecT{{\mathbf{T}}}
\def\vecN{{\mathbf{N}}}
\def\vecB{{\mathbf{B}}}
\def\vecv{{\mathbf{v}}}
\def\vecr{{\mathbf{r}}}
\def\vecf{{\mathbf{f}}}
\def\vecg{{\mathbf{g}}}
\def\vecupsilon{{\mathbf{\upsilon}}}
\def\vecu{{\mathbf{u}}}
\def\vecj{{\mathbf{j}}}
\def\vecx{{\mathbf{x}}}
\def\vecy{{\mathbf{y}}}
\def\vecz{{\mathbf{z}}}
\def\veck{{\mathbf{k}}}
\def\vecp{{\mathbf{p}}}
\def\vecn{{\mathbf{n}}}
\def\vecA{{\mathbf{A}}}
\def\vecP{{\mathbf{P}}}
\def\vecsigma{{\mathbf{\sigma}}}
\def\hatJn{{\hat{J_\vecn}}}
\def\hatJx{{\hat{J_x}}}
\def\hatJy{{\hat{J_y}}}
\def\hatJz{{\hat{J_z}}}
\def\hatj#1{\hat{J_{#1}}}
\def\hatphi{{\hat{\phi}}}
\def\hatq{{\hat{q}}}
\def\hatpi{{\hat{\pi}}}
\def\vel{\upsilon}
\def\Dint{{\mathcal{D}}}
\def\adag{{\hat{a}^\dagger}}
\def\bdag{{\hat{b}^\dagger}}
\def\cdag{{\hat{c}^\dagger}}
\def\ddag{{\hat{d}^\dagger}}
\def\hata{{\hat{a}}}
\def\hatb{{\hat{b}}}
\def\hatc{{\hat{c}}}
\def\hatd{{\hat{d}}}
\def\hatD{{\,\hat{D}}}
\def\hatN{{\hat{N}}}
\def\hatH{{\hat{H}}}
\def\hatp{{\hat{p}}}
\def\Fup{{F^{\mu\nu}}}
\def\Fdown{{F_{\mu\nu}}}
\def\newl{\nonumber \\}
\def\vece{\mathrm{e}}
\def\calM{{\mathcal{M}}}
\def\calT{{\mathcal{T}}}
\def\calR{{\mathcal{R}}}
\def\barpsi{\bar{\psi}}
\def\baru{\bar{u}}
\def\barv{\bar{\upsilon}}
\def\qeq{\stackrel{?}{=}}
\def\ftf{\stackrel{\mathcal{FT}}{\Longrightarrow}}
\def\ftb{\stackrel{\mathcal{FT}}{\Longleftarrow}}
\def\ftfb{\stackrel{\mathcal{FT}}{\Longleftrightarrow}}
\def\ltf{\stackrel{\mathcal{LT}}{\Longrightarrow}}
\def\ltb{\stackrel{\mathcal{LT}}{\Longleftarrow}}
\def\ltfb{\stackrel{\mathcal{LT}}{\Longleftrightarrow}}
\def\torder#1{\mathcal{T}\left(#1\right)}
\def\rorder#1{\mathcal{R}\left(#1\right)}
\def\contr#1#2{\contraction{}{#1}{}{#2}#1#2}
\def\trof#1{\mathrm{Tr}\left(#1\right)}
\def\trace{\mathrm{Tr}}
\def\comm#1{\ \ \ \left(\mathrm{used}\ #1\right)}
\def\tcomm#1{\ \ \ (\text{#1})}
\def\slp{\slashed{p}}
\def\slk{\slashed{k}}
\def\calp{{\mathfrak{p}}}
\def\veccalp{\mathbf{\mathfrak{p}}}
\def\Tthree{T_{\tiny \textcircled{3}}}
\def\pthree{p_{\tiny \textcircled{3}}}
\def\dbar{{\,\mathchar'26\mkern-12mu d}}
\def\erf{\mathrm{erf}}
\def\const{\mathrm{const.}}
\def\pheat{\pfrac p{\ln T}V}
\def\vheat{\pfrac V{\ln T}p}
%%units
\def\fdeg{{^\circ \mathrm{F}}}
\def\cdeg{^\circ \mathrm{C}}
\def\atm{\,\mathrm{atm}}
\def\angstrom{\,\text{\AA}}
\def\SIL{\,\mathrm{L}}
\def\SIkm{\,\mathrm{km}}
\def\SIyr{\,\mathrm{yr}}
\def\SIGyr{\,\mathrm{Gyr}}
\def\SIV{\,\mathrm{V}}
\def\SImV{\,\mathrm{mV}}
\def\SIeV{\,\mathrm{eV}}
\def\SIkeV{\,\mathrm{keV}}
\def\SIMeV{\,\mathrm{MeV}}
\def\SIGeV{\,\mathrm{GeV}}
\def\SIcal{\,\mathrm{cal}}
\def\SIkcal{\,\mathrm{kcal}}
\def\SImol{\,\mathrm{mol}}
\def\SIN{\,\mathrm{N}}
\def\SIHz{\,\mathrm{Hz}}
\def\SIm{\,\mathrm{m}}
\def\SIcm{\,\mathrm{cm}}
\def\SIfm{\,\mathrm{fm}}
\def\SImm{\,\mathrm{mm}}
\def\SInm{\,\mathrm{nm}}
\def\SImum{\,\mathrm{\mu m}}
\def\SIJ{\,\mathrm{J}}
\def\SIW{\,\mathrm{W}}
\def\SIkJ{\,\mathrm{kJ}}
\def\SIs{\,\mathrm{s}}
\def\SIkg{\,\mathrm{kg}}
\def\SIg{\,\mathrm{g}}
\def\SIK{\,\mathrm{K}}
\def\SImmHg{\,\mathrm{mmHg}}
\def\SIPa{\,\mathrm{Pa}}
%page
\def\secpage#1#2{\begin{frame}\bcenter{\bf \Huge #1} \ecenter \skipline \tbox{\bcenter #2 \ecenter}\end{frame}}
\def\append#1#2{\secpage{附录#1}{#2}}
\def\thinka#1{\begin{frame}\frametitle{思考题}\bcenter\lfig{0.4}{think0.jpg}  \skipline #1 \ecenter\end{frame}}
\def\thinkb#1{\begin{frame}\frametitle{思考题}\bcenter\lfig{0.5}{think1.jpg}  \skipline #1 \ecenter\end{frame}}
\def\thinkc#1{\begin{frame}\frametitle{思考题}\bcenter\lfig{0.5}{think2.jpg}  \skipline #1 \ecenter\end{frame}}
\def\thinkd#1{\begin{frame}\frametitle{思考题}\bcenter\lfig{0.5}{think3.jpg}  \skipline #1 \ecenter\end{frame}}
\def\thinke#1{\begin{frame}\frametitle{思考题}\bcenter\lfig{0.5}{think4.jpg}  \skipline #1 \ecenter\end{frame}}
\def\thinkf#1{\begin{frame}\frametitle{思考题}\bcenter\lfig{0.5}{think5.jpg}  \skipline #1 \ecenter\end{frame}}
\def\schw{ds^2 = \left(1-\frac{2GM}{r}\right)dt^2 - \left(1-\frac{2GM}{r}\right)^{-1} dr^2 - r^2\left(d\theta^2 + \sin^2\theta d\phi^2\right)}

\input{titlepage.tex}
  \date{}
  \begin{document}
  \bch
  \tpage{10}{From Riemann Tensor to Einstein Tensor}


  \begin{frame}
    \frametitle{该在哪里投影？}
    设 $P$, $Q$ 是很接近的两点，矢量 $\vec{A}$ 的协变微分给出的是 $\vec{A}(Q)$ 和 $\vec{A}(P)$ 按照 $P$ 处或者 $Q$ 处的投影规则的投影的差。


    \skiplines
    
    具体到底是按 $P$ 处还是 $Q$ 处的投影规则？我们曾评价说这不重要，因为差异是二阶小量。但是，在某些场景下，一阶小量完美抵消，二阶小量也会变得重要——

    \addfig{1}{catsaway.jpg}
  \end{frame}


  \begin{frame}
    \frametitle{二阶协变微商}
    在物理学中，除了热学中的相变等特殊情况之外，我们总是假定足够的光滑性，也就是高阶的普通微商可以任意交换微商次序。协变微商是否也是如此？

    \skiplines

    先说答案：No！ {\blue 高阶协变微商一般不能随意交换微商次序。}

    \skiplines
    
    问题还是出在投影的不一致性上。在考虑一阶微商时，我们通过取同一个投影规则解决了投影差异带来的（非物理）贡献。但是，这样有个二阶小量的不确定性——而二次微商考察的正是二阶小量。

  \end{frame}

  \begin{frame}

    \frametitle{想不清就硬怼……}
    由于标量场的第一次微商不涉及投影，标量场的二次微商其实可以交换次序。如果想不大清楚的话，可以来硬怼一下： 设 $S$ 为标量场

    $$ S_{;\mu;\nu} = S_{;\mu,\nu}-\Gamma^\alpha_{\ \mu\nu}S_{;\alpha}  = S_{,\mu,\nu} - \Gamma^\alpha_{\ \mu\nu}S_{,\alpha} $$
    

    推导竟出乎意料地简单——不过要记得感谢广义相对论用了对称的 Cristoffel 联络——或者说，感谢我们的时空那么单纯简洁。
  \end{frame}

  \begin{frame}
    \frametitle{这次运气就没那么好了}
    我们来计算一个矢量场 $A_\mu$ 的二阶微商：
    \bea
    && A_{\mu;\alpha;\beta} - A_{\mu;\beta;\alpha} \newl
    &=& A_{\mu;\alpha,\beta}-\Gamma^\rho_{\ \alpha\beta}A_{\mu;\rho} -\Gamma^\rho_{\ \mu\beta}A_{\rho;\alpha}- [\alpha\leftrightarrow\beta] \newl
    &=& A_{\mu;\alpha,\beta} -\Gamma^\rho_{\ \mu\beta}A_{\rho;\alpha}- [\alpha\leftrightarrow\beta] \newl    
    &=& \left(A_{\mu,\alpha}-\Gamma^\rho_{\ \mu\alpha}A_\rho\right)_{,\beta} -\Gamma^\rho_{\ \mu\beta}\left(A_{\rho,\alpha}-\Gamma^\lambda_{\ \rho\alpha}A_\lambda\right) - [\alpha\leftrightarrow\beta]\newl
    &=& -\Gamma^\rho_{\ \mu\alpha,\beta}A_\rho + \Gamma^\rho_{\ \mu\beta}\Gamma^\lambda_{\ \rho\alpha}A_\lambda- [\alpha\leftrightarrow\beta]\newl
    &=&\left(-\Gamma^\lambda_{\ \mu\alpha,\beta} + \Gamma^\rho_{\ \mu\beta}\Gamma^\lambda_{\ \rho\alpha}-[\alpha \leftrightarrow\beta]\right) A_\lambda\newl
    &\equiv & - R^\lambda_{\ \mu\alpha\beta}A_\lambda
    \eea
   最后一步只是定义。由于$A$是任意的，这里 $R^\lambda_{\ \mu\alpha\beta}$ 显然是张量。
  \end{frame}
  
  \begin{frame}
    \frametitle{黎曼张量(Riemann Tensor)和一些不愉快的符号差异}
    \tbox{$$R^\lambda_{\ \mu\alpha\beta} = \Gamma^\lambda_{\ \mu\alpha,\beta} + \Gamma^\rho_{\ \mu\alpha}\Gamma^\lambda_{\ \rho\beta} -[\alpha \leftrightarrow\beta]$$}
    这是一个让人不愉快的定义，因为如果你手头有很多相关的书，你会发现有些书是按上式定义的，而有些书则相差一个负号。

    \skipline
    
    这种不愉快还会延续到对里奇张量(Ricci Tensor)的定义，我们的定义是收缩第一个和最后一个指标：
    \tbox{$$ R_{\mu\nu} = R^\rho_{\ \mu\nu\rho}$$}
    而有些书上的定义也相差一个负号(收缩的是第一和第三个指标)。

  \end{frame}

  \begin{frame}
    \frametitle{物理学家找麻烦的决心}
    如果上述两次不愉快都发生了，那么两次相差的负号就互相抵消，最后本课程的核心方程——爱因斯坦的引力方程的符号就一致了。

    但如果你觉得所有人都会遵循让爱因斯坦方程一致的游戏规则，那你就太低估物理学家找麻烦的能力了——爱因斯坦方程同样有两种版本，一种是本课程采用的
    \tbox{$$R^{\mu\nu}-\frac{1}{2}R\,g^{\mu\nu}=8\pi GT^{\mu\nu}$$}
    另一种右边相差个负号。

    \skiplines

    在爱因斯坦方程里，左边的 $R=R^\alpha_{\ \alpha}$ 是里奇标量(Ricci scalar)，右边的 $G = 6.674 \times 10^{-11}\mathrm{N\,m^2/kg^2} $ 是牛顿万有引力常数， $T^{\mu\nu}$ 是物质的能量动量张量.
  \end{frame}
  
  \begin{frame}
    \frametitle{黎曼张量的协变形式}
    在讨论二维曲面时，我们曾介绍了协变形式的黎曼张量：
    \tbox{$$R_{\lambda\mu\alpha\beta} = \Gamma_{\lambda\mu\alpha,\beta}  + \Gamma_{\rho\lambda\alpha}\Gamma^\rho_{\ \mu\beta} - [\alpha\leftrightarrow \beta]$$}
    下面证明它和前面的定义是一致的：
    \bea
    R_{\lambda\mu\alpha\beta} &=&g_{\lambda\nu}R^\nu_{\ \mu\alpha\beta} \newl
    &=&  g_{\lambda\nu}\Gamma^\nu_{\ \mu\alpha,\beta} + g_{\lambda\nu} \Gamma^\rho_{\ \mu\alpha}\Gamma^\nu_{\ \rho\beta} -[\alpha \leftrightarrow\beta] \newl
    &=& \Gamma_{\lambda\mu\alpha,\beta} - g_{\lambda\nu,\beta}\Gamma^\nu_{\ \mu\alpha}+ g_{\lambda\nu} \Gamma^\rho_{\ \mu\alpha}\Gamma^\nu_{\ \rho\beta} -[\alpha \leftrightarrow\beta] \newl
    &=& \Gamma_{\lambda\mu\alpha,\beta} - \left(\Gamma_{\lambda\nu\beta}+\Gamma_{\nu\lambda\beta}\right)\Gamma^\nu_{\ \mu\alpha}+  \Gamma^\rho_{\ \mu\alpha}\Gamma_{\lambda\rho\beta} -[\alpha \leftrightarrow\beta] \newl
    &=& \Gamma_{\lambda\mu\alpha,\beta} - \Gamma_{\nu\lambda\beta}\Gamma^\nu_{\ \mu\alpha} -[\alpha \leftrightarrow\beta] \newl
    &=& \Gamma_{\lambda\mu\alpha,\beta} + \Gamma_{\nu\lambda\alpha}\Gamma^\nu_{\ \mu\beta} -[\alpha \leftrightarrow\beta]     
    \eea
  \end{frame}

  \begin{frame}
    下面我们讨论一个特殊的坐标系：

    \addfig{1}{taolunyiyi.jpg}
    
  \end{frame}
  
  \begin{frame}
    \frametitle{让一点的联络消失的魔术}
    如果在一个坐标系 $x^\mu$ 下某点 P 的联络为 $\Gamma^\lambda_{\ \mu\nu}(P)$，
    对在 P 的附近的任意一点 Q，令
    $$ \tilde{x}^\lambda(Q) = x^\lambda(Q) + \frac{1}{2}\Gamma^\lambda_{\ \mu\nu}(P) [x^\mu(Q)-x^\mu(P)][x^\nu(Q)-x^\nu(P)].$$
    该映射在 P 点满足 $\frac{\partial \tilde{x}^\mu}{\partial x^\nu} = \delta^\mu_{\ \nu}.$ 求 Jacobbi 矩阵的逆，还能得到 $\frac{\partial x^\mu}{\partial \tilde{x}^\nu} = \delta^\mu_{\ \nu}.$ 另外，
    $$\frac{\partial^2 \tilde{x}^\lambda}{\partial x^\mu\partial x^\nu} = \Gamma^\lambda_{\ \mu\nu}(P) .$$
    在P点的联络满足变换式
    $$ \Gamma^\lambda_{\ \mu\nu}(P) = \frac{\partial^2 \tilde{x}^\beta}{\partial x^\mu\partial x^\nu}\frac{\partial x^\lambda}{\partial\tilde{x}^\beta} + \widetilde{\Gamma}^\rho_{\ \alpha\beta}(P)\frac{\partial x^\lambda}{\partial\tilde{x}^\rho}\frac{\partial\tilde{x}^\alpha}{\partial x^\mu}\frac{\partial\tilde{x}^\beta}{\partial x^\nu} = \Gamma^\lambda_{\ \mu\nu}(P) + \widetilde{\Gamma}^\lambda_{\ \mu\nu}(P).$$
    也就是在新的 $\tilde{x}^\mu$ 坐标系里，P 点的联络消失了！
  \end{frame}

  \begin{frame}
    \frametitle{在魔术坐标系里的黎曼张量}
    在魔术坐标系里，由于所有的联络在 P 点消失，P 点的黎曼张量满足：
    \bea
    && 2R_{\lambda\mu\alpha\beta}\newl
    &=& 2\Gamma_{\lambda\mu\alpha,\beta} -  2\Gamma_{\lambda\mu\beta,\alpha}\newl
    &=& g_{\lambda\mu,\alpha,\beta} + g_{\lambda\alpha,\mu,\beta}-g_{\mu\alpha,\lambda,\beta}-g_{\lambda\mu,\beta,\alpha}-g_{\lambda\beta,\mu,\alpha}+g_{\mu\beta,\lambda,\alpha} \newl
    &=&  g_{\lambda\alpha,\mu,\beta}-g_{\mu\alpha,\lambda,\beta}-g_{\lambda\beta,\mu,\alpha}+g_{\mu\beta,\lambda,\alpha}     
    \eea
    于是容易验证:
  \end{frame}

  \begin{frame}
    \frametitle{黎曼张量的对称性}
    \tbox{$$R_{\lambda\mu\alpha\beta}=R_{\alpha\beta\lambda\mu};$$
      $$R_{\lambda\mu\alpha\beta}=-R_{\lambda\mu\beta\alpha};$$
      $$R_{\lambda\mu\alpha\beta}=-R_{\mu\lambda\beta\alpha};$$
      $$R_{\lambda\mu\alpha\beta}+R_{\lambda\alpha\beta\mu}+R_{\lambda\beta\mu\alpha}=0.$$
    }        
    由于张量方程在坐标变换下不变，所以这些对称性在任何坐标系成立。
  \end{frame}

  \begin{frame}
    在魔术坐标系里到底发生了什么呢？我们仍然以嵌入三维欧氏空间的曲面为例，用 $\vecx$ 来标记三维欧氏空间的矢量。联络
    $$\Gamma_{\lambda\mu\nu} = \vecx_{,\lambda}\cdot\vecx_{,\mu,\nu}$$
    全部消失的意思就是所有 $\vecx_{,\mu,\nu}$ 都和曲面上的所有切向量正交。

    
    也就是说，P 点所有 $\vecx_{,\mu,\nu}$ 或者是零， 或者沿曲面的法向量 $\vecn$ 的方向。{\blue 如果曲面上在 P 点的爬虫采用了这个坐标系，它无法发现空间是弯曲的（切向量的变化都在爬虫看不到的维度里）。}

    \addfig{1}{jiangwei2.jpg}
    
    这为我们后面讨论的等效原理埋下了伏笔。
  \end{frame}


  \begin{frame}
    \frametitle{在魔术坐标系更容易看出来对称性}
    在魔术坐标系里 P 点的黎曼张量
    $$ R_{\lambda\mu\alpha\beta}= \Gamma_{\lambda\mu\alpha,\beta} -  \Gamma_{\lambda\mu\beta,\alpha} = \vecx_{,\lambda,\beta}\cdot\vecx_{,\mu,\alpha} - \vecx_{,\lambda,\alpha}\cdot\vecx_{,\mu,\beta}.$$
    注意所有 $\vecx_{,\rho,\sigma}$ 都可以替换成 $ \vecn\cdot \vecx_{,\rho,\sigma}$，上式把曲面的第二基本形式和黎曼张量联系了起来。从这个表达式也更容易看出黎曼张量的各种对称性。
  \end{frame}

  \begin{frame}
    在魔术坐标系里的P点
    $$2R_{\lambda\mu\alpha\beta}  =g_{\lambda\alpha,\mu,\beta}  -g_{\mu\alpha,\lambda,\beta}-g_{\lambda\beta,\mu,\alpha}+g_{\mu\beta,\lambda,\alpha} $$      
    P点联络消失，可以尽情地用普通微商替代协变微商
    $$ 2R_{\lambda\mu\alpha\beta;\gamma} = g_{\lambda\alpha,\mu,\beta,\gamma}  -g_{\mu\alpha,\lambda,\beta,\gamma}-g_{\lambda\beta,\mu,\alpha,\gamma}+g_{\mu\beta,\lambda,\alpha,\gamma}$$
    $$ 2R_{\lambda\mu\beta\gamma;\alpha} = g_{\lambda\beta,\mu,\gamma,\alpha}  -g_{\mu\beta,\lambda,\gamma,\alpha}-g_{\lambda\gamma,\mu,\beta,\alpha}+g_{\mu\gamma,\lambda,\beta,\alpha}$$    
    $$ 2R_{\lambda\mu\gamma\alpha;\beta} = g_{\lambda\gamma,\mu,\alpha,\beta}  -g_{\mu\gamma,\lambda,\alpha,\beta}-g_{\lambda\alpha,\mu,\gamma,\beta}+g_{\mu\alpha,\lambda,\gamma,\beta}$$    
    把这三个式子加起来，就得到
  \end{frame}


  \begin{frame}
    \frametitle{Bianchi恒等式}
    \tbox{$$ R_{\lambda\mu\alpha\beta;\gamma} + R_{\lambda\mu\beta\gamma;\alpha} + R_{\lambda\mu\gamma\alpha;\beta} = 0. $$}
    由于这是个张量方程，所以在一般坐标系也成立。
    
    把 $\lambda$ 和 $\beta$ 缩并，以及 $\mu$ 和 $\alpha$ 缩并，得到：
    $$ R_{;\gamma} - 2R^\alpha_{\ \gamma;\alpha} = 0.$$
    由于度规的协变微商为零，上式可以写成：
    $$ \left(R^\alpha_{\ \gamma}-\frac{1}{2}g^\alpha_{\ \gamma}R\right)_{;\alpha} = 0.$$    
  \end{frame}


  \begin{frame}
    \frametitle{爱因斯坦张量}
    或者写成更广为人知的形式
    {\blue $$ G^{\mu\nu}_{\ \ ;\mu} = 0.$$}    
    这里的“爱因斯坦张量”定义为
    {\blue $$ G^{\mu\nu} \equiv R^{\mu\nu}-\frac{1}{2}g^{\mu\nu}R$$}
    是一个对称张量。

    由于物质能量动量张量守恒 $T^{\mu\nu}_{\ \ ;\mu}=0$，爱因斯坦大胆假设 $G^{\mu\nu}$ 和 $T^{\mu\nu}$ 成正比，得到了著名的爱因斯坦方程。我们下一讲对它进行详细介绍。
  \end{frame}
  
\ech
\end{document}




